\documentclass[10pt,a4paper]{amsart}
\usepackage[margin=19mm]{geometry}
\usepackage{amsmath,amssymb,mathtools}
\usepackage[scr=rsfs,cal=euler]{mathalfa}
\usepackage{xfrac}
\usepackage[unicode,hidelinks,bookmarks=false]{hyperref}
\newcommand{\ie}{i.e.\ }
\newcommand{\cf}{cf.\ }
\newcommand{\resp}{respectively\ }
\newcommand{\Subsection}{\subsection}
\providecommand{\nobreakdash}{\nobreak\hbox{-}}
\setlength{\emergencystretch}{2em}
\multlinegap0pt
\usepackage{amssymb}
\usepackage{bm}
\usepackage{algorithm}\usepackage{algpseudocode}
\usepackage{mathtools}
\newtheorem{lemma}{Lemma}
  \def\alpha{alpha}%
\newcommand{\mmse}{\mathrm{mmse}}
\title{When Diffusion Model Can Ignore Dimension: An Entropy-Based Theory}
\author{Ahmad Aghapour, Erhan Bayraktar}
\date{Source: arXiv:2605.07969v1 (CC BY 4.0)}
\begin{document}
\maketitle

\noindent\emph{Source and licence.} When Diffusion Model Can Ignore Dimension: An Entropy-Based Theory. Version arXiv:2605.07969v1 (CC BY 4.0), distributed by arXiv under the Creative Commons Attribution 4.0 International licence, http://creativecommons.org/licenses/by/4.0/, as recorded in the arXiv licence metadata for this exact version and shown on the abstract page https://arxiv.org/abs/2605.07969v1. Full licence notice: \texttt{licenses/arxiv-2605.07969-CC-BY-4.0.txt}.\par\smallskip

\noindent\textit{Editorial scope.} Counted unit: the article's lemma lem:hybrid-grid-bound (source0.tex lines 1105--1132). The article's complete original proof of it is delivered with it (source0.tex lines 1134--1381, one \texttt{proof} environment of 248 lines). No mathematics is written, repaired or re-derived: the statement and the proof are the article's own lines, in the article's own notation, and no retained line is substituted or re-flowed.  No further source-local context is needed: every cross-reference of every retained line resolves inside the two delivered spans.  Nothing was omitted from any retained span, so the package carries no omission record.  No retained line contains a citation, so the wrapper carries no bibliography and no bibliography entry.  No retained line contains a figure environment, an image-inclusion command or drawing code, so no image is delivered and the package declares no asset dependency.  Editorial formatting only, in addition: an AMS \texttt{amsart} preamble that restates the article's own declarations and macros used by the retained lines, namely the environments \texttt{lemma} on separate counters, together with the standard packages those lines need.  The retained environments print in this document as Lemma 1 in order of appearance, whereas the article prints the same items with the numbering of its own full document.  The complete native source of the article as distributed in the arXiv e-print for this exact version is bundled unchanged as \texttt{sources/200123/source0.tex}.
\begin{lemma}
\label{lem:hybrid-grid-bound}
Fix $h\in(0,\alpha]$. There exists a grid
\[
    \eta_{\min}=\eta_0<\eta_1<\cdots<\eta_N=\eta_{\max}
\]
such that
\[
    N
    \le
    \left\lceil \frac{(\alpha-\eta_{\min})_+}{h}\right\rceil
    +
    \left\lceil
    \frac{\log_+(\eta_{\max}/\alpha)}
    {\log(1+h/\alpha)}
    \right\rceil
    +2
\]
and
\[
    \mathcal E_{\mathrm{disc}}
    \le
    hR
    \left(
    2+\log_+\frac{\eta_{\max}}{\alpha}
    \right).
\]
\end{lemma}

\begin{proof}
We give a precise construction of the grid and then bound the discretization error by summation by parts. To avoid notational distractions, we first treat the main case
\[
    \eta_{\min}<\alpha<\eta_{\max}.
\]
The cases $\alpha\le \eta_{\min}$ and $\eta_{\max}\le\alpha$ are obtained by removing, respectively, the uniform or the geometric part of the construction.

Fix $h\in(0,\alpha]$ and set
\[
    r:=1+\frac{h}{\alpha}.
\]
Let
\[
    M:=\left\lceil \frac{\alpha-\eta_{\min}}{h}\right\rceil .
\]
We define the uniform part of the grid by
\[
    \eta_i:=\eta_{\min}+ih,
    \qquad
    0\le i\le M-1,
\]
and then set
\[
    \eta_M:=\alpha.
\]
Thus the first $M-1$ increments are equal to $h$, while the last increment in the uniform part is
\[
    \eta_M-\eta_{M-1}
    =
    \alpha-\eta_{\min}-(M-1)h
    \in(0,h].
\]
This last interval may be shorter than $h$, and we will account for this explicitly below.

For the geometric part, let
\[
    Q:=\left\lceil
    \frac{\log(\eta_{\max}/\alpha)}{\log r}
    \right\rceil .
\]
Starting from $\eta_M=\alpha$, define
\[
    \eta_{M+j}:=\alpha r^j,
    \qquad
    1\le j\le Q-1,
\]
and finally set
\[
    \eta_{M+Q}:=\eta_{\max}.
\]
Since $Q$ is the first integer for which $\alpha r^Q\ge \eta_{\max}$, the final geometric interval is possibly shortened, and all previous geometric intervals have the exact multiplicative ratio $r$. The total number of intervals is
\[
    N=M+Q .
\]
Moreover,
\[
    M\le \frac{\alpha-\eta_{\min}}{h}+1
    \le \frac{\alpha}{h}+1
\]
and
\[
    Q\le
    \frac{\log(\eta_{\max}/\alpha)}{\log r}+1 .
\]
Using $\log(1+x)\ge x/(1+x)$ with $x=h/\alpha$, and using $h\le \alpha$, we obtain
\[
    \log r
    =
    \log\left(1+\frac{h}{\alpha}\right)
    \ge
    \frac{h}{\alpha+h}
    \ge
    \frac{h}{2\alpha}.
\]
Hence
\[
    Q
    \le
    \frac{2\alpha}{h}
    \log\frac{\eta_{\max}}{\alpha}
    +1 .
\]
Therefore
\[
    N
    \le
    \frac{\alpha}{h}
    \left(
    1+2\log\frac{\eta_{\max}}{\alpha}
    \right)
    +2 .
\]

We now bound the discretization error. Write
\[
    m_k:=\mmse(\eta_k),
    \qquad
    \Delta_k:=\eta_k-\eta_{k-1}.
\]
By the MMSE representation,
\[
    \mathcal E_{\mathrm{disc}}
    =
    \sum_{k=1}^N
    \int_{\eta_{k-1}}^{\eta_k}
    \bigl(m_{k-1}-\mmse(\eta)\bigr)\,d\eta .
\]
Since $\mmse$ is nonincreasing,
\[
    \mmse(\eta)\ge m_k,
    \qquad
    \eta\in[\eta_{k-1},\eta_k],
\]
and therefore
\[
    \mathcal E_{\mathrm{disc}}
    \le
    \sum_{k=1}^N
    \Delta_k(m_{k-1}-m_k).
\]
The summation-by-parts identity gives
\[
\begin{aligned}
    \sum_{k=1}^N
    \Delta_k(m_{k-1}-m_k)
    &=
    \Delta_1m_0
    +
    \sum_{k=1}^{N-1}
    (\Delta_{k+1}-\Delta_k)m_k
    -
    \Delta_Nm_N .
\end{aligned}
\]
Since the last term is nonpositive, we discard it.

We now control the remaining terms carefully. On the uniform part, the increments are equal to $h$ except possibly for the last interval ending at $\alpha$. Hence the interior differences are zero until the shortened interval. The initial contribution is bounded by
\[
    \Delta_1m_0\le hR,
\]
because $\Delta_1\le h$ and $m_0\le R$.

There is one possible positive jump at the transition from the last uniform interval to the first geometric interval. Indeed, the first geometric interval has length
\[
    \alpha r-\alpha=\alpha(r-1)=h,
\]
while the last uniform interval has length at most $h$. Therefore this transition contributes at most
\[
    h\,\mmse(\alpha)\le hR.
\]
This is the term that must be included explicitly.

It remains to control the increases inside the geometric part. For the nontruncated geometric intervals,
\[
    \eta_{M+j}=\alpha r^j,
    \qquad
    j\ge 0,
\]
and the corresponding increments are
\[
    \Delta_{M+j+1}
    =
    \eta_{M+j+1}-\eta_{M+j}
    =
    \alpha r^j(r-1).
\]
Thus, for consecutive full geometric intervals,
\[
    \Delta_{M+j+2}-\Delta_{M+j+1}
    =
    \alpha r^j(r-1)^2 .
\]
At the endpoint $\eta_{M+j+1}=\alpha r^{j+1}$, the entropy bound on the MMSE gives
\[
    m_{M+j+1}
    \le
    \frac{2H(J)}{\alpha r^{j+1}} .
\]
Therefore each positive increase inside the geometric region is bounded by
\[
\begin{aligned}
    \bigl(\Delta_{M+j+2}-\Delta_{M+j+1}\bigr)m_{M+j+1}
    &\le
    \alpha r^j(r-1)^2
    \frac{2H(J)}{\alpha r^{j+1}}  \\
    &=
    \frac{2H(J)(r-1)^2}{r}.
\end{aligned}
\]
The final geometric interval may be shortened. A shortened final interval cannot create a larger positive increase than the corresponding full geometric interval, so the same bound applies to the last possible increase as well.

The number of positive increases inside the geometric part is at most $Q-1$. Since
\[
    Q-1
    \le
    \frac{\log(\eta_{\max}/\alpha)}{\log r},
\]
the total geometric contribution is at most
\[
    \frac{2H(J)(r-1)^2}{r}
    \frac{\log(\eta_{\max}/\alpha)}{\log r}.
\]
Using
\[
    \log r\ge \frac{r-1}{r},
\]
we get
\[
    \frac{(r-1)^2}{r\log r}
    \le r-1.
\]
Hence the geometric contribution is bounded by
\[
    2H(J)(r-1)
    \log\frac{\eta_{\max}}{\alpha}.
\]
Since $r-1=h/\alpha$ and $\alpha=2H(J)/R$, this becomes
\[
    hR\log\frac{\eta_{\max}}{\alpha}.
\]

Combining the initial uniform contribution, the transition contribution, and the geometric contribution yields
\[
    \mathcal E_{\mathrm{disc}}
    \le
    hR
    \left(
    2+\log\frac{\eta_{\max}}{\alpha}
    \right).
\]
This proves the desired bound in the case $\eta_{\min}<\alpha<\eta_{\max}$.

If $\eta_{\max}\le\alpha$, the geometric part is absent. The same argument gives
\[
    \mathcal E_{\mathrm{disc}}\le hR,
\]
which is stronger than the displayed bound with the logarithmic term removed. If $\alpha\le\eta_{\min}$, the uniform part is absent and the same geometric argument gives
\[
    \mathcal E_{\mathrm{disc}}
    \le
    hR
    \log\frac{\eta_{\max}}{\eta_{\min}}
    \le
    hR
    \log\frac{\eta_{\max}}{\alpha},
\]
using $\eta_{\min}\ge\alpha$. Thus, in all cases, the hybrid-grid estimate holds after replacing the logarithm by $\log_+(\eta_{\max}/\alpha)$ and allowing the universal constant in front of the nonlogarithmic term.
\end{proof}

\end{document}
